\setcounter{section}{22}

Sekarang kita membahas suatu teknik bantu analitik penting yang disebut \stichwort {partisi satuan} {.} Kita akan menggunakannya dalam pembuktian pernyataan bahwa manifold yang dapat diorientasikan mempunyai bentuk volume positif, dan dalam pembuktian teorema Stokes. Dalam kuliah ini kita akan mengonstruksi partisi satuan; untuk itu mula-mula kita memerlukan beberapa konsep topologis.

  \zwischenueberschrift{Penghabisan kompak}

\bild{ \begin{center}
\includegraphics[width=5.5cm]{ \bildeinlesungpng {Inner_point} {png} }
\end{center}
\bildtext {Interior terbuka adalah gabungan semua titik interior, yaitu titik-titik dari $T$
\zusatzklammer {pada gambar, $E$} {} {,}
yang memiliki suatu lingkungan terbuka yang seluruhnya termuat di dalam $T$.} }

\bildlizenz { Inner point.png } {} {Zasdfgbnm} {Commons} {PD} {}

\inputdefinition
{}
{

Misalkan $X$ suatu
\definitionsverweis {ruang topologis}{}{}
dan

\mavergleichskette
{\vergleichskette
{T
}
{ \subseteq }{ X
}
{  }{
}
{    }{
}
{   }{
}
}
{}{}{}
suatu himpunan bagian. Maka

\mavergleichskettedisp
{\vergleichskette
{ \overline{  T  }
}
{ =} {  \bigcap_{A \text{ tertutup}, \, T \subseteq A} A
}
{ } {
}
{ } {
}
{ } {
}
}
{}{}{}
disebut \definitionswort {penutupan}{}
\zusatzklammer {atau \definitionswort {penutupan topologis}{}} {} {}
dari $T$.

}
Untuk ruang metrik, kita sebelumnya telah memperkenalkan penutupan sebagai himpunan semua
\definitionsverweis {titik pelekatan}{}{}{;}
lihat khususnya
Soal 35.3 (Analisis (Osnabrück 2021--2023)).

\inputdefinition
{}
{

Misalkan $X$ suatu
\definitionsverweis {ruang topologis}{}{}
dan

\mavergleichskette
{\vergleichskette
{ T
}
{ \subseteq }{ X
}
{  }{
}
{    }{
}
{   }{
}
}
{}{}{}
suatu himpunan bagian. Maka

\mavergleichskettedisp
{\vergleichskette
{ T ^{o}
}
{ =} { \bigcup_{U \text{ terbuka}, \, U \subseteq T}   U
}
{ } {
}
{ } {
}
{ } {
}
}
{}{}{}
disebut \definitionswort {interior terbuka}{}
\zusatzklammer {atau \stichwort {interior} {}} {} {}
dari $T$.

}
Kedua konsep ini berhubungan melalui

\mavergleichskettedisp
{\vergleichskette
{ \overline{ T }
}
{ =} { X \setminus (X \setminus T)^{o}
}
{ } {
}
{ } {
}
{ } {
}
}
{}{}{.}
Konsep batas juga terbawa dari situasi metrik ke ruang topologis sembarang.

\inputdefinition
{
}
{

Misalkan $X$ suatu
\definitionsverweis {ruang topologis}{}{}
dan

\mavergleichskette
{\vergleichskette
{ T
}
{ \subseteq }{ X
}
{  }{
}
{    }{
}
{   }{
}
}
{}{}{}
suatu himpunan bagian. Yang dimaksud dengan
\definitionswort {batas}{}
dari $T$ adalah himpunan

\mavergleichskettedisp
{\vergleichskette
{ \operatorname{Batas} \, ( T )
}
{ =} { \overline{  T  }  \setminus T ^{o}
}
{ } {
}
{ } {
}
{ } {
}
}
{}{}{.}

}
Perhatikan bahwa batas topologis ini merupakan konsep yang berbeda dari batas suatu manifold berbatas. Meskipun demikian, terdapat hubungan yang dibahas dalam
Soal 22.1.

\inputdefinition
{}
{

Misalkan $X$ suatu
\definitionsverweis {ruang topologis}{}{}
dan
\maabbdisp {f} { X } {\R
} {}
suatu
\definitionsverweis {fungsi}{}{.}
Maka
\definitionsverweis {penutupan topologis}{}{}

\mathdisp {\overline{  \left\{  x \in X  \mid  f(x) \neq 0        \right\}   }} {  }

disebut \definitionswort {tumpuan}{} dari $f$.

}

\inputdefinition
{}
{

Misalkan $X$ suatu
\definitionsverweis {ruang topologis}{}{.}
Suatu \definitionswort {penghabisan kompak}{}

\mathbed {A_n} {}
 {n \in \N} {}
 {} {} {} {,}
dari $X$ adalah suatu
\definitionsverweis {barisan}{}{}
\definitionsverweis {himpunan-himpunan bagian kompak}{}{}

\mavergleichskette
{\vergleichskette
{ A_n
}
{ \subseteq }{ X
}
{  }{
}
{    }{
}
{   }{
}
}
{}{}{}
dengan

\mathdisp {A_n \subseteq A_{n+1}^{o} \text{ dan } \bigcup_{n=0}^\infty A_n = X} { . }

}

Ruang $\R^d$ dihabiskan secara kompak oleh bola-bola tertutup

\mathbed {B  \left( 0,n \right)} {}
 {n \in \N} {}
 {} {} {} {.}
Lihat
Soal 22.6.

  \zwischenueberschrift{Partisi satuan}



\inputfaktbeweis
{Mannigfaltigkeit/Abzählbare Basis/Kompakte Ausschöpfung/Fakt}
{Lema}
{}
{

\faktsituation {Misalkan $M$ suatu
\definitionsverweis {manifold}{}{}}
\faktvoraussetzung {dengan suatu
\definitionsverweis {basis terhitung bagi topologinya}{}{.}}
\faktfolgerung {Maka $M$ mempunyai suatu
\definitionsverweis {penghabisan kompak}{}{.}}
\faktzusatz {}
\faktzusatz {}

}
{

Untuk setiap titik

\mavergleichskette
{\vergleichskette
{ P
}
{ \in }{ M
}
{  }{
}
{    }{
}
{   }{
}
}
{}{}{}
terdapat suatu lingkungan koordinat terbuka

\mavergleichskette
{\vergleichskette
{ P
}
{ \in }{ U
}
{  }{
}
{    }{
}
{   }{
}
}
{}{}{,}
\maabbdisp {\alpha} { U } { V
} {}
serta lingkungan-lingkungan bola

\mavergleichskettedisp
{\vergleichskette
{ U \left(  \alpha(P), \epsilon  \right)
}
{ \subseteq} { B  \left( \alpha(P), \epsilon \right)
}
{ \subset} { V
}
{ } {
}
{ } {
}
}
{}{}{.}
Karena pemetaan koordinat tersebut merupakan
\definitionsverweis {homeomorfisme}{}{}
dan bola tertutup bersifat kompak, maka

\mavergleichskette
{\vergleichskette
{ B_P
}
{ = }{ \alpha^{-1} \left(  B  \left( \alpha(P), \epsilon \right)  \right)
}
{  }{
}
{    }{
}
{   }{
}
}
{}{}{}
merupakan suatu
\definitionsverweis {himpunan bagian kompak}{}{}
dari $M$ yang memuat
\definitionsverweis {lingkungan terbuka}{}{}

\mavergleichskette
{\vergleichskette
{ U_P
}
{ = }{ \alpha^{-1} \left(  U \left(  \alpha(P), \epsilon  \right)  \right)
}
{  }{
}
{    }{
}
{   }{
}
}
{}{}{}
dari $P$. Himpunan-himpunan

\mathbed {U_P} {}
 {P \in M} {}
 {} {} {} {,}
membentuk suatu
\definitionsverweis {tutupan terbuka}{}{}
dari $M$, sehingga menurut
Soal 2.8 (Teori Ukuran dan Integrasi (Osnabrück 2022--2023))
terdapat suatu subtutupan terhitung. Nyatakan subtutupan ini dengan

\mathbed {U_n} {}
 {n \in \N} {}
 {} {} {} {,}
\zusatzklammer {dengan $U_n$ termuat di dalam himpunan kompak $B_n$} {} {.}
Sekarang kita definisikan secara
\definitionsverweis {rekursif}{}{}
suatu pemetaan
\definitionsverweis {naik secara monoton}{}{}
\maabbeledisp {} {\N} {\N
} { k } {n_k
} {,}
sedemikian sehingga

\mathbeddisp {A_k = \bigcup_{n=0}^{n_k } B_n} {}
 {k \in \N} {}
 {} {} {} {,}
merupakan suatu
\definitionsverweis {penghabisan kompak}{}{}
dari $M$. Sebagai gabungan berhingga dari himpunan-himpunan kompak, setiap $A_k$ bersifat kompak. Kita mulai dengan

\mavergleichskette
{\vergleichskette
{ n_0
}
{ = }{ 0
}
{  }{
}
{    }{}
{   }{}
}
{}{}{.}
Misalkan $n_k$ sudah dikonstruksi. Himpunan

\mathdisp {A_k \cup B_{n_k +1}} {  }

bersifat
\definitionsverweis {kompak}{}{}
dan karena itu ditutupi oleh berhingga banyak himpunan terbuka
\mathl{\bigcup_{n=0}^{n_{k+1} } U_n}{,} dengan memilih

\mavergleichskette
{\vergleichskette
{ n_{k+1}
}
{ \geq }{ n_k +1
}
{  }{
}
{    }{
}
{   }{
}
}
{}{}{.}
Dengan pilihan ini,

\mavergleichskettedisp
{\vergleichskette
{ A_k
}
{ \subseteq} { \bigcup_{n = 0}^{n_{k+1} } U_n
}
{ \subseteq} { \bigcup_{n = 0}^{n_{k+1} } B_n
}
{ =} { A_{k+1}
}
{ } {
}
}
{}{}{,}
dan barisan ini membentuk suatu penghabisan karena himpunan-himpunan

\mathbed {U_n} {}
 {n \in \N} {}
 {} {} {} {,}
membentuk suatu tutupan terbuka dari $M$.

}

\bild{ \begin{center}
\includegraphics[width=5.5cm]{ \bildeinlesungsvg {Partition_of_unity_illustration} {svg} }
\end{center}
\bildtext {Ilustrasi beberapa fungsi taknegatif yang jumlahnya selalu 1.} }

\bildlizenz { Partition of unity illustration.svg } {} {Oleg Alexandrov} {Commons} {PD} {}

\inputdefinition
{}
{

Misalkan $X$ suatu
\definitionsverweis {ruang topologis}{}{.}
Suatu keluarga fungsi
\maabbdisp {h_j} { X } {\R
} {}
dengan

\mavergleichskette
{\vergleichskette
{ j
}
{ \in }{ J
}
{  }{
}
{    }{
}
{   }{
}
}
{}{}{}
disebut suatu \definitionswort {partisi satuan}{,} jika sifat-sifat berikut berlaku.
\aufzaehlungdrei{Berlaku

\mavergleichskette
{\vergleichskette
{ h_j(X)
}
{ \subseteq }{ [0,1]
}
{  }{
}
{    }{
}
{   }{
}
}
{}{}{}
untuk setiap

\mavergleichskette
{\vergleichskette
{ j
}
{ \in }{ J
}
{  }{
}
{    }{
}
{   }{
}
}
{}{}{.}
}{Setiap titik

\mavergleichskette
{\vergleichskette
{ P
}
{ \in }{ X
}
{  }{
}
{    }{
}
{   }{
}
}
{}{}{}
mempunyai suatu
\definitionsverweis {lingkungan terbuka}{}{}

\mavergleichskette
{\vergleichskette
{ P
}
{ \in }{ U
}
{  }{
}
{    }{
}
{   }{
}
}
{}{}{}
sedemikian sehingga
\definitionsverweis {fungsi-fungsi hasil pembatasan}{}{}

\mathl{h_j {{|}}_{ U }}{} semuanya, kecuali berhingga banyak, merupakan fungsi nol.
}{Berlaku

\mavergleichskette
{\vergleichskette
{  \sum_{j \in J} h_j
}
{ = }{ 1
}
{  }{
}
{    }{
}
{   }{
}
}
{}{}{.}
} Jika semua $h_j$ \definitionsverweis {kontinu}{}{,} partisi tersebut disebut suatu \definitionswort {partisi satuan kontinu}{.}

}

Sifat kedua memastikan bahwa jumlah dalam (3) terdefinisi, sebab untuk setiap titik

\mavergleichskette
{\vergleichskette
{ P
}
{ \in }{ X
}
{  }{
}
{    }{
}
{   }{
}
}
{}{}{}
dan bagi semua indeks

\mavergleichskette
{\vergleichskette
{ j
}
{ \in }{ J
}
{  }{
}
{    }{
}
{   }{
}
}
{}{}{}
kecuali berhingga banyak, berlaku kesamaan

\mavergleichskette
{\vergleichskette
{  h_j(P)
}
{ = }{ 0
}
{  }{
}
{    }{
}
{   }{
}
}
{}{}{.}
Pada suatu manifold, partisi seperti itu disebut \stichwort {diferensiabel} {,} jika semua $h_j$ merupakan
\definitionsverweis {fungsi-fungsi diferensiabel}{}{.}

\inputdefinition
{}
{

Misalkan

\mavergleichskette
{\vergleichskette
{ X
}
{ = }{ \bigcup_{i \in I} W_i
}
{  }{
}
{    }{
}
{   }{
}
}
{}{}{}
suatu
\definitionsverweis {tutupan terbuka}{}{}
dari suatu
\definitionsverweis {ruang topologis}{}{}
$X$. Suatu
\definitionsverweis {partisi satuan}{}{}
\maabbdisp {h_j} { X } {\R
} {}
dengan

\mavergleichskette
{\vergleichskette
{ j
}
{ \in }{ J
}
{  }{
}
{    }{
}
{   }{
}
}
{}{}{}
disebut suatu \definitionswort {partisi satuan subordinat terhadap tutupan}{,} jika untuk setiap

\mavergleichskette
{\vergleichskette
{ j
}
{ \in }{ J
}
{  }{
}
{    }{
}
{   }{
}
}
{}{}{}
terdapat suatu himpunan terbuka
\mathl{W_{i(j)}}{} dari tutupan tersebut sedemikian sehingga
\definitionsverweis {tumpuan}{}{}
dari $h_j$ termuat di dalam
\mathl{W_{i(j)}}{}{.}

}



\inputfaktbeweis
{Mannigfaltigkeit/Abzählbare Basis/Lokal endlicher Atlas/Fakt}
{Lema}
{}
{

\faktsituation {Misalkan $M$ suatu
\definitionsverweis {manifold}{}{}
dengan suatu
\definitionsverweis {basis terhitung}{}{}
bagi topologinya. Misalkan

\mavergleichskette
{\vergleichskette
{ M
}
{ = }{ \bigcup_{i \in I} W _i
}
{  }{
}
{    }{
}
{   }{
}
}
{}{}{}
suatu
\definitionsverweis {tutupan terbuka}{}{}
dari $M$.}
\faktfolgerung {Maka terdapat suatu
\definitionsverweis {atlas kompatibel}{}{}
yang terhitung

\mathbed {\left(  U_j, \alpha_j, V_j  \right)} {}
 {j \in J} {}
 {} {} {} {,}
dengan lingkungan-lingkungan bola

\mavergleichskettedisp
{\vergleichskette
{ U \left(   0 , \delta_j   \right)
}
{ \subset} { B  \left(  0 , \epsilon_j  \right)
}
{ \subset} { V_j
}
{ } {
}
{ } {
}
}
{}{}{}
\zusatzklammer {di sini,

\mavergleichskettek
{\vergleichskettek
{ 0
}
{ \in }{ V_j
}
{ \subseteq }{ \R^n
}
{    }{
}
{   }{
}
}
{}{}{}
dan

\mavergleichskettek
{\vergleichskettek
{ \delta_j
}
{ < }{ \epsilon_j
}
{  }{
}
{    }{
}
{   }{
}
}
{}{}{}} {} {}
sedemikian sehingga untuk setiap

\mavergleichskette
{\vergleichskette
{ j
}
{ \in }{ J
}
{  }{
}
{    }{
}
{   }{
}
}
{}{}{}
terdapat suatu
\mathl{W_{i(j)}}{} dengan

\mavergleichskette
{\vergleichskette
{ U_j
}
{ \subseteq }{ W_{i(j)}
}
{  }{
}
{    }{
}
{   }{
}
}
{}{}{}; selain itu, $M$ ditutupi oleh

\mathbed {\alpha_j^{-1} \left(  U \left(   0 , \delta_j   \right)  \right)} {}
 {j \in J} {}
 {} {} {} {,}
dan setiap titik

\mavergleichskette
{\vergleichskette
{ P
}
{ \in }{ M
}
{  }{
}
{    }{
}
{   }{
}
}
{}{}{}
terletak hanya di dalam berhingga banyak himpunan $U_j$.}
\faktzusatz {}
\faktzusatz {}

}
{

\teilbeweis {}{}{}
{Misalkan tutupan terbuka

\mathbed {W_i} {}
 {i \in I} {}
 {} {} {} {,}
diberikan. Selanjutnya, misalkan

\mathbed {A_n} {}
 {n \in \N} {}
 {} {} {} {,}
suatu
\definitionsverweis {penghabisan kompak}{}{}
dari $M$, yang keberadaannya dijamin oleh
Lema 22.6.
Himpunan-himpunan terbuka
\mathl{A_{n+2}^{o}  \setminus A_{n-1}}{} juga membentuk suatu tutupan terbuka, sebab untuk setiap titik

\mavergleichskette
{\vergleichskette
{ P
}
{ \in }{ M
}
{  }{
}
{    }{
}
{   }{
}
}
{}{}{}
terdapat suatu nilai minimal

\mavergleichskette
{\vergleichskette
{ n
}
{ \in }{ \N
}
{  }{
}
{    }{
}
{   }{
}
}
{}{}{}
dengan

\mavergleichskette
{\vergleichskette
{ P
}
{ \in }{ A_n
}
{  }{
}
{    }{
}
{   }{
}
}
{}{}{}
\zusatzklammer {tetapkan

\mavergleichskettek
{\vergleichskettek
{ A_{-1}
}
{ = }{ \emptyset
}
{  }{
}
{    }{
}
{   }{
}
}
{}{}{}} {} {.}
Untuk $n$ ini berlaku
\mathkor {} {P \not\in A_{n-1}} {dan} {P \in A_n \subseteq A_{n+1} ^{o}} {.}
Dengan meninjau irisan-irisan
\mathl{W_i \cap \left(  A_{n+2}^{o}  \setminus A_{n-1}  \right)}{}{,} kita dapat menganggap bahwa setiap himpunan dari tutupan termuat di dalam suatu
\mathl{A_{n+2}^{o} \setminus A_{n-1}}{}{.}}
{}
\teilbeweis {}{}{}
{Untuk setiap titik

\mavergleichskette
{\vergleichskette
{ P
}
{ \in }{ M
}
{  }{
}
{    }{
}
{   }{
}
}
{}{}{}
terdapat suatu lingkungan koordinat terbuka
\zusatzklammer {yang kompatibel} {} {}

\mavergleichskette
{\vergleichskette
{ P
}
{ \in }{ U_P
}
{  }{
}
{    }{
}
{   }{
}
}
{}{}{,}
yang termuat di dalam salah satu $W_i$ dan yang mempunyai lingkungan-lingkungan bola

\mavergleichskettedisp
{\vergleichskette
{ U \left(   0 , \delta_P   \right)
}
{ \subset} { B  \left(  0 , \epsilon_P  \right)
}
{ \subset} { V_P
}
{ } {
}
{ } {
}
}
{}{}{}
dengan

\mavergleichskette
{\vergleichskette
{ P
}
{ \in }{ \alpha_P^{-1} \left(  U \left(   0 , \delta_P   \right)  \right)
}
{  }{
}
{    }{
}
{   }{
}
}
{}{}{}
dan

\mavergleichskette
{\vergleichskette
{ \delta_P
}
{ < }{ \epsilon_P
}
{  }{
}
{    }{
}
{   }{
}
}
{}{}{.}
Himpunan-himpunan

\mathbed {\alpha_P^{-1} \left(  U \left(   0 , \delta_P   \right)  \right)} {}
 {P \in M} {}
 {} {} {} {,}
kemudian juga membentuk suatu tutupan terbuka dari $M$. Menurut
Soal 2.8 (Teori Ukuran dan Integrasi (Osnabrück 2022--2023))
kita dapat beralih ke suatu subtutupan terhitung. Jadi, kita dapat menganggap bahwa suatu sistem domain peta

\mathbed {U_j} {}
 {j \in \N} {}
 {} {} {} {,}
bersama dengan lingkungan-lingkungan bola

\mavergleichskettedisp
{\vergleichskette
{ U \left(   0 , \delta_j   \right)
}
{ \subset} { B  \left(  0 , \epsilon_j  \right)
}
{ \subset} { V_j
}
{ } {
}
{ } {
}
}
{}{}{}
diberikan sedemikian sehingga

\mathbed {\alpha_j^{-1} \left(  U \left(   0 , \delta_j   \right)  \right)} {}
 {j \in \N} {}
 {} {} {} {,}
juga merupakan suatu tutupan terbuka dari $M$, setiap $U_j$ termuat di dalam suatu
\mathl{W_{i(j)}}{}{,} dan hubungan dengan penghabisan kompak yang diuraikan di atas berlaku.}
{}
\teilbeweis {}{Kita akan mendefinisikan suatu himpunan bagian

\mavergleichskette
{\vergleichskette
{ J
}
{ \subseteq }{ \N
}
{  }{
}
{    }{
}
{   }{
}
}
{}{}{}
sedemikian sehingga keluarga

\mathbed {U_j} {}
 {j \in J} {}
 {} {} {} {,}
juga memenuhi sifat keberhinggaan tersebut.\leerzeichen{}}{}
{Mula-mula, kekompakan $A_0$ memberikan suatu himpunan indeks berhingga
\mathl{J_{-1}}{} sedemikian sehingga himpunan-himpunan

\mathbed {\alpha_j^{-1} \left(  U \left(   0 , \delta_j   \right)  \right)} {}
 {j \in J_{-1}} {}
 {} {} {} {,}
menutupi $A_0$. Untuk

\mavergleichskette
{\vergleichskette
{ n
}
{ \in }{ \N
}
{  }{
}
{    }{
}
{   }{
}
}
{}{}{}
tinjau himpunan kompak
\mathl{A_{n+1} \setminus A_n ^{o}}{.} Himpunan ini ditutupi oleh berhingga banyak himpunan

\mathbed {\alpha_j^{-1} \left(  U \left(   0 , \delta_j   \right)  \right)} {}
 {j \in \N} {}
 {} {} {} {,}
dan untuk itu hanya diperlukan indeks-indeks $j$ sedemikian sehingga $U_j$ termuat di dalam
\mathl{A_{n+2}^{o}  \setminus A_{n-1}}{}{.} Nyatakan himpunan indeks berhingga yang bersangkutan dengan $J_n$, dan tetapkan

\mavergleichskette
{\vergleichskette
{ J
}
{ = }{ J_{-1} \cup \bigcup_{n \in \N} J_n
}
{  }{
}
{    }{
}
{   }{
}
}
{}{}{.}
Maka setiap $A_n$ hanya memotong berhingga banyak himpunan

\mathbed {U_j} {}
 {j \in J} {}
 {} {} {} {.}}
{}

}




\inputfaktbeweis
{Mannigfaltigkeit/Abzählbare Basis/Überdeckung/Partition/Fakt}
{Teorema}
{}
{

\faktsituation {Misalkan $M$ suatu
\definitionsverweis {manifold diferensiabel}{}{}
dengan suatu
\definitionsverweis {basis terhitung}{}{}
bagi topologinya.}
\faktfolgerung {Maka untuk setiap tutupan terbuka terdapat suatu
\definitionsverweis {partisi satuan}{}{}
yang \definitionsverweis {terdiferensialkan secara kontinu}{}{}
dan subordinat terhadap tutupan tersebut.}
\faktzusatz {}
\faktzusatz {}

}
{

Menurut
Lema 22.9
kita dapat menganggap bahwa suatu tutupan terbuka terdiri atas domain-domain peta

\mathbed {U_j} {}
 {j \in J} {}
 {} {} {} {,}
\zusatzklammer {$J$ terhitung} {} {}
dengan
\maabbdisp {\alpha_j} {U_j } { V_j
} {}
dan dengan lingkungan-lingkungan bola

\mavergleichskettedisp
{\vergleichskette
{ U \left(   0 , \delta_j   \right)
}
{ \subset} { B  \left(  0 , \epsilon_j  \right)
}
{ \subset} { V_j
}
{ } {
}
{ } {
}
}
{}{}{}
\zusatzklammer {dengan

\mavergleichskettek
{\vergleichskettek
{ \delta_j
}
{ < }{ \epsilon_j
}
{  }{
}
{    }{
}
{   }{
}
}
{}{}{}} {} {}
sedemikian sehingga himpunan-himpunan
\mathl{\alpha_j^{-1} \left(  U \left(   0 , \delta_j   \right)  \right)}{} juga membentuk suatu tutupan dari $M$ dan setiap titik

\mavergleichskette
{\vergleichskette
{ P
}
{ \in }{ M
}
{  }{
}
{    }{
}
{   }{
}
}
{}{}{}
termuat hanya dalam berhingga banyak himpunan $U_j$ dan, khususnya, hanya dalam berhingga banyak himpunan
\mathl{\alpha_j^{-1} \left(  U \left(   0 , \delta_j   \right)  \right)}{}{.}
Pada $V_j$, tinjau fungsi $g_j$ yang didefinisikan oleh

\mavergleichskettedisp
{\vergleichskette
{ g_j (v)
}
{ =} { \begin{cases}  e^{ - { \frac{   1  }{  \left(  \delta_j^2- \Vert { v } \Vert^2  \right)^2 } } }  \text{ untuk } \Vert { v } \Vert < \delta_j \, , \\ 0 \text{ selainnya} \, , \end{cases}
}
{ } {
}
{ } {
}
{ } {
}
}
{}{}{}
Fungsi ini bernilai positif tepat pada
\mathl{U \left(   0 , \delta_j   \right)}{} dan
\definitionsverweis {tumpuan}{}{}
fungsi ini adalah
\mathl{B  \left(  0 , \delta_j  \right)}{.}
Peninjauan pada kedua himpunan bagian terbuka
\zusatzklammer {yang menutupi $V_j$} {} {}
\mathkor {} {U \left(   0 , \epsilon_j   \right)} {dan} {V_j \setminus B  \left(  0 , \delta_j  \right)} {}
menunjukkan bahwa $g_j$ terdiferensialkan tak hingga kali. Kita mendefinisikan suatu fungsi
\maabbdisp {\tilde{g}_j} { M } {\R
} {}
melalui

\mavergleichskettedisp
{\vergleichskette
{ \tilde{g}_j(x)
}
{ =} { \begin{cases}  g_j( \alpha_j(x)) \text{ untuk }  x \in \alpha_j^{-1}( U \left(   0 , \delta_j   \right) ) \, , \\  0 \text{ selainnya} \, . \end{cases}
}
{ } {
}
{ } {
}
{ } {
}
}
{}{}{}
Fungsi ini
\definitionsverweis {terdiferensialkan secara kontinu}{}{}
pada $M$, sebab \anfuehrung{pita}{}
\mathl{B  \left(  0 , \epsilon_j  \right) \setminus U \left(   0 , \delta_j   \right)}{} memungkinkan suatu transisi mulus. Kita tetapkan

\mavergleichskettedisp
{\vergleichskette
{ \tilde{g}(x)
}
{ \defeq} { \sum_{j \in J} \tilde{g}_j(x)
}
{ } {
}
{ } {
}
{ } {
}
}
{}{}{,}
yang merupakan jumlah berhingga untuk setiap titik, sebab
\definitionsverweis {tumpuan}{}{}
dari $\tilde{g}_j$ termuat di dalam

\mavergleichskettedisp
{\vergleichskette
{ \alpha_j^{-1} \left(  B  \left(  0 , \delta_j  \right)  \right)
}
{ \subset} { U_j
}
{ } {
}
{ } {
}
{ } {
}
}
{}{}{.}
Fungsi ini
\definitionsverweis {terdiferensialkan secara kontinu}{}{}
pada $M$ dan positif di setiap titik, sebab
\mathl{\tilde{g}_j(x)}{} masing-masing bernilai positif pada himpunan yang bersesuaian dari tutupan
\mathl{\alpha_j^{-1} \left(  U \left(   0 , \delta_j   \right)  \right)}{}{.}
Dengan demikian, fungsi-fungsi

\mavergleichskettedisp
{\vergleichskette
{ h_j
}
{ =} { { \frac{   \tilde{g}_j  }{ \tilde{g}  } }
}
{ } {
}
{ } {
}
{ } {
}
}
{}{}{}
membentuk partisi satuan yang dicari.

}

  \zwischenueberschrift{Orientasi pada manifold dan bentuk volume}

Dengan bantuan partisi satuan, sekarang kita dapat membuktikan kebalikan dari
Lema 15.5.



\inputfaktbeweis
{Nullstellenfreie Volumenform/Orientierte (Karten) Mannigfaltigkeit/Äquivalenz/Fakt}
{Teorema}
{}
{

\faktsituation {Misalkan $M$ suatu
\definitionsverweis {manifold diferensiabel
}{}{}
dengan suatu
\definitionsverweis {basis terhitung bagi topologinya}{}{.}}
\faktfolgerung {Maka terdapat suatu
\definitionsverweis {bentuk volume}{}{}
\definitionsverweis {kontinu}{}{}
tanpa titik nol pada $M$ jika dan hanya jika $M$
\definitionsverweis {dapat diorientasikan}{}{.}}
\faktzusatz {Bentuk volume ini juga dapat dipilih agar terdiferensialkan secara kontinu dan positif.}
\faktzusatz {}

}
{

\teilbeweis {}{}{}
{Salah satu arah sudah dibuktikan
dalam Lema 15.5.}
{}
\teilbeweis {}{}{}
{Sebaliknya, misalkan $M$
\definitionsverweis {dapat diorientasikan}{}{}
dan diberikan suatu
\definitionsverweis {atlas berorientasi}{}{}
yang
\definitionsverweis {terhitung}{}{}

\mathbed {\left(  U_i,\alpha_i,V_i  \right)} {}
 {i \in I} {}
 {} {} {} {,}
dari $M$. Di sini,

\mavergleichskette
{\vergleichskette
{ V_i
}
{ \subseteq }{ \R^n
}
{  }{
}
{    }{
}
{   }{
}
}
{}{}{}
merupakan himpunan
\definitionsverweis {terbuka}{}{,}
dan
\definitionsverweis {koordinat-koordinat}{}{}

\mathl{x_1 , \ldots , x_n}{} mendefinisikan suatu
\definitionsverweis {bentuk volume}{}{}
kontinu tanpa titik nol
\zusatzklammer {bahkan terdiferensialkan tak hingga kali} {} {}

\mathl{dx_1 \wedge \ldots \wedge dx_n}{} pada $V_i$. Kita tetapkan

\mavergleichskettedisp
{\vergleichskette
{ \omega_i
}
{ =} { \alpha_i^* \left( dx_1 \wedge \ldots \wedge dx_n \right)
}
{ } {
}
{ } {
}
{ } {
}
}
{}{}{}
dan dengan demikian memperoleh suatu bentuk volume tanpa titik nol pada $U_i$, yang kita perluas dengan $0$ di luar $U_i$\zusatzfussnote {Perluasan ini tentu saja tidak kontinu, tetapi hal itu tidak menjadi masalah untuk pembahasan berikut} {.} {.}

Sekarang, misalkan

\mathbed {h_j} {}
 {j \in J} {}
 {} {} {} {,}
suatu keluarga fungsi yang subordinat terhadap tutupan

\mathbed {U_i} {}
 {i \in I} {}
 {} {} {} {,}
dan terdiferensialkan secara kontinu, serta membentuk suatu
\definitionsverweis {partisi satuan}{}{,}
yang keberadaannya dijamin oleh
Teorema 22.10.
Jadi, khususnya, untuk setiap $j$ terdapat suatu
\mathl{i(j)}{} sedemikian sehingga
\definitionsverweis {tumpuan}{}{}
dari $h_j$ termuat di dalam
\mathl{U_{i(j)}}{}{.}
Oleh karena itu,
\mathl{h_j \omega_{i(j)}}{} merupakan suatu
\definitionsverweis {bentuk diferensial}{}{}
kontinu berderajat $n$ pada $M$. Kita tetapkan

\mavergleichskettedisp
{\vergleichskette
{ \omega
}
{ =} { \sum_{j \in J} h_j \omega_{i(j)}
}
{ } {
}
{ } {
}
{ } {
}
}
{}{}{.}
Di setiap titik

\mavergleichskette
{\vergleichskette
{ P
}
{ \in }{ M
}
{  }{
}
{    }{
}
{   }{
}
}
{}{}{}
jumlah ini berhingga dan karena itu mendefinisikan suatu bentuk diferensial kontinu berderajat $n$ pada $M$. Untuk suatu titik

\mavergleichskette
{\vergleichskette
{ P
}
{ \in }{ M
}
{  }{
}
{    }{
}
{   }{
}
}
{}{}{}
dan suatu
\definitionsverweis {basis}{}{}

\mathl{v_1 , \ldots , v_n}{} dari
\mathl{T_PM}{} yang merepresentasikan orientasi, berlaku

\mavergleichskettedisp
{\vergleichskette
{ \omega(P; v_1 , \ldots , v_n)
}
{ =} { \sum_{j \in J} h_j(P) \omega_{i(j)} \left(  P; v_1 , \ldots , v_n  \right)
}
{ } {
}
{ } {
}
{ } {
}
}
{}{}{.}
Terdapat suatu $j$ dengan

\mavergleichskette
{\vergleichskette
{ h_j(P)
}
{ > }{ 0
}
{  }{
}
{    }{
}
{   }{
}
}
{}{}{,}
dan untuk $j$ ini juga berlaku

\mavergleichskette
{\vergleichskette
{ \omega_{i(j)}(P; v_1 , \ldots , v_n)
}
{ > }{ 0
}
{  }{
}
{    }{
}
{   }{
}
}
{}{}{,}
sebab

\mavergleichskette
{\vergleichskette
{ P
}
{ \in }{ U_{i(j)}
}
{  }{
}
{    }{
}
{   }{
}
}
{}{}{}
berlaku; dengan demikian, bentuk ini positif di setiap titik.}
{}

}
